\documentclass[12pt]{article}

% --- Packages ---
\usepackage[utf8]{inputenc}
\usepackage[T1]{fontenc}
\usepackage[french]{babel}
\usepackage{amsmath, amssymb}
\usepackage{geometry}
\usepackage{tikz}
\usepackage[most]{tcolorbox}
\usepackage{lmodern}
\usepackage{pgfplots}
\pgfplotsset{compat=1.18}
\usepackage{mdframed}
\usepackage{bm}
\usepackage{xcolor}
\usepackage{sectsty}
\usetikzlibrary{calc}

\usepackage[utf8]{inputenc}
\usepackage[T1]{fontenc}
\usepackage[dvipsnames]{xcolor}
\usepackage{tkz-tab}
\usepackage{multido}
% --- Page ---
\geometry{margin=2cm}

% Footer
\usepackage{fancyhdr}
\pagestyle{fancy}
\cfoot{\thepage}
\rfoot{\textit{G.M.}}
\renewcommand{\headrulewidth}{0pt}
\renewcommand{\footrulewidth}{0pt}

% Pour les normes : 
\newcommand*{\norme}[1]{\left\lVert\overrightarrow{#1}\right\rVert}

\usetikzlibrary{shapes.misc}

\tikzset{cross/.style={cross out, draw=black, minimum size=2*(#1-\pgflinewidth), inner sep=0pt, outer sep=0pt},
%default radius will be 1pt. 
cross/.default={4pt}}

% mettre dans le préambule, après \usepackage{xcolor}
\usepackage{sectsty} % gère facilement la couleur des titres
\setlength{\parindent}{0pt}

\sectionfont{\color{red}\Large\bfseries}     % pour \section*
% \subsectionfont{\color{red}\normalsize\bfseries} % pour \subsection*

\begin{document}

% ============================
%   VECTEURS
% ============================

% \begin{center}
%     {\LARGE \textbf{Vecteurs}}\\[4mm]
%     {\large 2nde générale}
% \end{center}

% \hrule
% \vspace{0.4cm}

% ============================
%   PARTIE 3 : Notions de vecteurs
% ============================

\section*{\textcolor{red}{Partie 3 : Produit d'un vecteur par un réel}}

\vspace{5pt}

\underline{\textbf{Exemple 1 :}}

\begin{minipage}[t]{.45\textwidth}%
    \vspace{5pt}
    $5\overrightarrow{u}$ est la somme de 5 vecteurs $\overrightarrow{u}$
    \vspace{5pt}\\On a :
    \vspace{5pt}\\$5\overrightarrow{u}$ = $\overrightarrow{u} + \overrightarrow{u} + \overrightarrow{u}
    + \overrightarrow{u} + \overrightarrow{u}$
\end{minipage}%
\hfill
\begin{minipage}[t]{.45\textwidth}%
    \begin{center}
        \begin{tikzpicture}[scale=1]

        % --- Vecteur u seul ---
        \draw[->,very thick,blue] (0,0) -- (1,0.6);
        \node[below right] at (0.3,0.3) {$\overrightarrow{u}$};

        % --- 5 vecteurs u mis bout à bout (à côté) ---
        \draw[->,very thick,blue] (2,0) -- (3,0.6);
        \draw[->,very thick,blue] (3,0.6) -- (4,1.2);
        \draw[->,very thick,blue] (4,1.2) -- (5,1.8);
        \draw[->,very thick,blue] (5,1.8) -- (6,2.4);
        \draw[->,very thick,blue] (6,2.4) -- (7,3);

        % \node[below right] at (4.5,1.2) {$5 \times \vec u$};

        % --- Vecteur 5u ---
        \draw[->,very thick,red] (3,-0.5) -- (8,2.5);
        \node[below right] at (5.5,1) {$5\overrightarrow{u}$};

        \end{tikzpicture}

    \end{center}
\end{minipage}%

\vspace{10pt}

\underline{Remarques :} 
\vspace{5pt}
\begin{itemize}
    \item Les vecteurs $5\overrightarrow{u}$ et 
    $\overrightarrow{u}$ ont la même direction et le même sens.
    \item La norme du vecteur $5\overrightarrow{u}$ est égale à 5 fois la norme du 
    vecteur $\overrightarrow{u}$.
\end{itemize}

\vspace{20pt}

\underline{\textbf{Exemple 2 :}}

\begin{minipage}[t]{.45\textwidth}%
    \vspace{5pt}
    $-3\overrightarrow{u}$ est la somme de 3 vecteurs $-\overrightarrow{u}$
    \vspace{5pt}\\On a :
    \vspace{5pt}\\$-3\overrightarrow{u}$ = $(-\overrightarrow{u}) + (-\overrightarrow{u}) 
    + (-\overrightarrow{u})$
\end{minipage}%
\hfill
\begin{minipage}[t]{.45\textwidth}%
    \vspace{0pt}
    \begin{center}
        \begin{tikzpicture}[scale=1]

        % --- Vecteur u seul ---
        \draw[->,very thick,blue] (0,0) -- (1,0.6);
        \node[below right] at (0.3,0.3) {$\overrightarrow{u}$};

        % --- 5 vecteurs u mis bout à bout (à côté) ---
        \draw[<-,very thick,blue] (2,0) -- (3,0.6);
        \draw[<-,very thick,blue] (3,0.6) -- (4,1.2);
        \draw[<-,very thick,blue] (4,1.2) -- (5,1.8);

        \node[below right] at (2.7,1.5) {$-\overrightarrow{u}$};
        \node[below right] at (1.7,0.9) {$-\overrightarrow{u}$};
        \node[below right] at (3.7,2.1) {$-\overrightarrow{u}$};

        % --- Vecteur 5u ---
        \draw[<-,very thick,red] (3,-0.5) -- (6,1.3);
        \node[below right] at (4.4,0.5) {$-3\overrightarrow{u}$};

        \end{tikzpicture}

    \end{center}
\end{minipage}%

\vspace{10pt}

\underline{Remarques :} 
\vspace{5pt}
\begin{itemize}
    \item Les vecteurs $\overrightarrow{u}$ et 
    $-\overrightarrow{u}$ ont la même direction et même norme mais de \textbf{sens contraire}.
    \item La norme du vecteur $-3\overrightarrow{u}$ est égale à 3 fois la norme du 
    vecteur $\overrightarrow{u}$.
\end{itemize}

\vspace{10pt}

% =============================================
% ================ EXERCICE 1 =================
% =============================================

\begin{mdframed}[
    linewidth=1pt,
    roundcorner=5pt,
    linecolor=RoyalBlue,
    innertopmargin=10pt,
    innerbottommargin=10pt
]
{\color{RoyalBlue}\textbf{\underline{Exercice d'application n°1 :} Représenter un vecteur
défini comme un produit et somme de vecteurs.}}
\vspace{10pt}

a)\vspace{5pt}\\
\begin{minipage}{.55\textwidth}%
    \begin{tikzpicture}[scale=0.5]

    % Grille seule
    \draw[step=1cm, gray, thin] (0,0) grid (17,8);

    % Vecteurs
    \draw[->,very thick,blue] (7,1) -- (14,1);
    \node[below right, blue] at (9.9,2.2) {$\overrightarrow{v}$};
    \draw[->,very thick,red] (7,1) -- (9,4);
    \node[below right, red] at (6.7,3.2) {$\overrightarrow{u}$};

    \end{tikzpicture}
\end{minipage}%
\hfill
\begin{minipage}{.45\textwidth}%
    Soient deux vecteurs $\overrightarrow{u}$ et $\overrightarrow{v}$. Représenter les 
    vecteurs suivants :
    \begin{center}
        $2\overrightarrow{u}$
        \vspace{5pt}\\
        $-\overrightarrow{v}$
        \vspace{5pt}\\
        $2\overrightarrow{u} - \overrightarrow{v}$
    \end{center}
\end{minipage}%

\vspace{10pt}
\underline{Méthode :}
\vspace{5pt}
\begin{itemize}
    \item On commence par représenter le vecteur $2\overrightarrow{u}$ : on place 
    bout-à-bout \textbf{deux vecteurs $\bm{\overrightarrow{u}}$}.
    \item Le vecteur $-\overrightarrow{v}$ a la même direction et la même longueur que 
    $\overrightarrow{v}$ mais il est de sens contraire.
    \item Pour représenter le vecteur $2\overrightarrow{u} - \overrightarrow{v} = 
    2\overrightarrow{u}+(-\overrightarrow{v})$, on place bout-à-bout les vecteurs 
    $2\overrightarrow{u}$ et $-\overrightarrow{v}$ et on relit les extrémités du chemin 
    construit.
\end{itemize}

\vspace{20pt}

b)\vspace{5pt}\\
\begin{minipage}{.45\textwidth}%
    \begin{tikzpicture}[scale=0.5]

    % Grille seule
    \draw[step=1cm, gray, thin] (0,0) grid (13,12);

    % Points
    \coordinate (A) at (7,7);
    \coordinate (B) at (7,9);
    \coordinate (C) at (9,10);

    % Points
    \fill (A) circle (3pt) node[below right] {$A$};
    \fill (B) circle (3pt) node[above left] {$B$};
    \fill (C) circle (3pt) node[above right] {$C$};

    \end{tikzpicture}
\end{minipage}%
\hfill
\begin{minipage}{.55\textwidth}%
    Soient 3 points $A$, $B$ et $C$.
    \vspace{5pt}\\Représenter le vecteur $\overrightarrow{BC} - 3\overrightarrow{AC}$.
\end{minipage}%

\vspace{10pt}

\underline{Méthode :}
\vspace{5pt}
\begin{itemize}
    \item Pour représenter le vecteur $\overrightarrow{BC} - 3\overrightarrow{AC} = 
    \overrightarrow{BC} + (-3\overrightarrow{AC})$, on place bout-à-bout les vecteurs 
    $\overrightarrow{BC}$ et $-3\overrightarrow{AC}$
\end{itemize}

\end{mdframed}

% =============================================
% ================ EXERCICE 2 =================
% =============================================

\begin{mdframed}[
    linewidth=1pt,
    roundcorner=5pt,
    linecolor=RoyalBlue,
    innertopmargin=10pt,
    innerbottommargin=10pt
]
{\color{RoyalBlue}\textbf{\underline{Exercice d'application n°2 :} Construire un point 
vérifiant une égalité vectorielle.}}
\vspace{10pt}

a)\vspace{5pt}\\
\begin{minipage}{.45\textwidth}%
    \begin{tikzpicture}[scale=0.5]

    % Grille seule
    \draw[step=1cm, gray, thin] (0,0) grid (13,12);

    % Point
    \coordinate (O) at (3,7);

    % Vecteurs
    \draw[->,very thick,blue] (7,11) -- (8,9);
    \node[below right, blue] at (6.1,10.5) {$\overrightarrow{u}$};
    \draw[->,very thick,red] (12,11) -- (9,10);
    \node[below right, red] at (10,10.5) {$\overrightarrow{v}$};

    \fill (O) circle (3pt) node[above left] {$O$};

    \end{tikzpicture}
\end{minipage}%
\hfill
\begin{minipage}{.55\textwidth}%
    Soient deux vecteurs $\overrightarrow{u}$ et $\overrightarrow{v}$ et un point $O$.
    \vspace{5pt}\\Construire le point $A$ tel que $\overrightarrow{OA} = 3\overrightarrow{u} - 
    \overrightarrow{v}$.
\end{minipage}%

\vspace{10pt}
\underline{Méthode :}
\vspace{5pt}
\begin{itemize}
    \item Pour représeter le vecteur $\overrightarrow{OA} = 3\overrightarrow{u} - \overrightarrow{v}$, 
    on place bout-à-bout les vecteurs $3\overrightarrow{u}$ et $-\overrightarrow{v}$ en partant 
    de $O$.
    \item Le point $A$ se trouve à l'extrémité du vecteur $-\overrightarrow{v}$ dans le chemin 
    construit.
\end{itemize}

\newpage

b)\vspace{5pt}\\
\begin{minipage}{.45\textwidth}%
    \begin{tikzpicture}[scale=0.5]

    % Grille seule
    \draw[step=1cm, gray, thin] (0,0) grid (13,12);

    % Points
    \coordinate (A) at (9,6);
    \coordinate (B) at (9,10);
    \coordinate (C) at (7,7);

    % Points
    \fill (A) circle (3pt) node[below right] {$A$};
    \fill (B) circle (3pt) node[above left] {$B$};
    \fill (C) circle (3pt) node[above right] {$C$};

    \end{tikzpicture}
\end{minipage}%
\hfill
\begin{minipage}{.55\textwidth}%
    Soient 3 points $A$, $B$ et $C$ du plan.
    \vspace{5pt}\\Construire le point $M$ tel que $\overrightarrow{AM} = -\overrightarrow{AB} + 
    3\overrightarrow{AC}$.
\end{minipage}%

\vspace{10pt}

\underline{Méthode :}
\vspace{5pt}
\begin{itemize}
    \item Pour représenter le vecteur $\overrightarrow{AM} = -\overrightarrow{AB} + 
    3\overrightarrow{AC}$, on place bout-à-bout les vecteurs 
    $-\overrightarrow{AB}$ et $3\overrightarrow{AC}$ en partant de $A$.
    \item Le point $M$ se trouve à l'extrémité du vecteur $3\overrightarrow{AC}$ dans le chemin 
    construit.
\end{itemize}

\end{mdframed}

% =============================================
% ================ EXERCICE 3 =================
% =============================================

\begin{mdframed}[
    linewidth=1pt,
    roundcorner=5pt,
    linecolor=RoyalBlue,
    innertopmargin=10pt,
    innerbottommargin=10pt
]
{\color{RoyalBlue}\textbf{\underline{Exercice d'application n°3 :} Exprimer par une lecture 
graphique un vecteur en fonction d'autres vecteurs.}}
\vspace{10pt}

\begin{minipage}{.35\textwidth}%
    \begin{tikzpicture}[scale=0.5]

    % Grille seule
    \draw[step=1cm, gray, thin] (0,0) grid (9,5);

    % Vecteurs
    \draw[->,very thick,blue] (2,1) -- (4,1);
    \node[below right, blue] at (2.2,1) {$\overrightarrow{a}$};

    \draw[->,very thick,red] (2,1) -- (2,2);
    \node[below right, red] at (0.5,2) {$\overrightarrow{b}$};

    \draw[->,very thick,black] (2,1) -- (8,4);
    \node[below right, black] at (5,2.7) {$\overrightarrow{u}$};

    \end{tikzpicture}
\end{minipage}%
\hfill
\begin{minipage}{.65\textwidth}%
    Par lecture graphique, exprimer le vecteur $\overrightarrow{u}$ en fonction des
    vecteurs $\overrightarrow{a}$ et $\overrightarrow{b}$.
\end{minipage}%

\vspace{10pt}
\underline{Méthode :}
\vspace{5pt}
\begin{itemize}
    \item On construit un chemin formé de vecteurs $\overrightarrow{a}$ et $\overrightarrow{b}$ 
    mis bout-à-bout reliant l’origine 
    et l’extrémité du vecteur $\overrightarrow{u}$.
    \item On compte ainsi le nombre de vecteurs $\overrightarrow{a}$ et $\overrightarrow{b}$ 
    formant ce chemin.
\end{itemize}

\end{mdframed}

\newpage

% ============================
%   PARTIE 4 : Notions de vecteurs
% ============================

\section*{\textcolor{red}{Partie 4 : Notion de colinéarité}}

\vspace{5pt}

\underline{\textbf{Exemple :}}

\begin{minipage}[t]{.40\textwidth}%
    \vspace{5pt}
    Les vecteurs $\overrightarrow{u}$ et $\overrightarrow{v}$ ont la \textbf{même direction}, 
    on dit qu'ils sont
    \textbf{colinéaires}.
\end{minipage}%
\hfill
\begin{minipage}[t]{.55\textwidth}%
    \vspace{0pt}
    \begin{center}
        \begin{tikzpicture}[scale=1]

        % --- Vecteur u seul ---
        \draw[->,very thick,blue] (0,0) -- (2,2);
        \node[below right] at (1,1) {$\overrightarrow{u}$};

        % --- 5 vecteurs u mis bout à bout (à côté) ---
        % \draw[->,very thick,blue] (2,0) -- (3,0.6);
        % \draw[->,very thick,blue] (3,0.6) -- (4,1.2);
        % \draw[->,very thick,blue] (4,1.2) -- (5,1.8);
        % \draw[->,very thick,blue] (5,1.8) -- (6,2.4);
        % \draw[->,very thick,blue] (6,2.4) -- (7,3);

        % \node[below right] at (4.5,1.2) {$5 \times \vec u$};

        % --- Vecteur 5u ---
        \draw[<-,very thick,red] (2,-1) -- (6,3);
        \node[below right] at (4,1) {$\overrightarrow{v} = -2\overrightarrow{u}$};

        \end{tikzpicture}

    \end{center}
\end{minipage}%

\vspace{10pt}

\begin{mdframed}[
    linewidth=1pt,
    roundcorner=5pt,
    linecolor=ForestGreen,
    innerbottommargin=10pt,
    innertopmargin=10pt
]
{\color{ForestGreen}\textbf{Définition :}}
Deux vecteurs non nuls $\overrightarrow{u}$ et $\overrightarrow{v}$ sont \textbf{colinéaires} signifie 
qu'ils ont la même direction (les vecteurs sont « portés » par des droites parallèles).
\vspace{5pt}\\C'est-à-dire qu'il existe un nombre réel $k$ tel que :
\[
\overrightarrow{u} = k\overrightarrow{v}
\]
\end{mdframed}

\vspace{10pt}

\underline{Remarque :} Le vecteur nul $\overrightarrow{0}$ est colinéaire à tout vecteur du plan.

\vspace{10pt}

\begin{mdframed}[
    linewidth=1pt,
    roundcorner=5pt,
    linecolor=RoyalBlue,
    innertopmargin=10pt,
    innerbottommargin=10pt
]
{\color{RoyalBlue}\textbf{\underline{Exercice d'application n°4 :} Démontrer que des vecteurs 
sont colinéaires.}}
\vspace{10pt}

On donne deux vecteurs $\overrightarrow{u}$ et $\overrightarrow{v}$, tel que :
\[
-4\overrightarrow{u} + 3\overrightarrow{v} = \overrightarrow{0}
\]

Démontrer que les vecteurs $\overrightarrow{u}$ et $\overrightarrow{v}$ sont colinéaires.

\end{mdframed}


\end{document}
